% TOP_PROBLEM: {"id": 30006782, "title": "Kimura–O'Sullivan Finite-Dimensionality Conjecture", "category_id": 5, "category": {"id": 5, "name": "algebraic_geometry", "display_name": "Algebraic Geometry", "description": "Geometric objects defined by polynomial equations.", "slug": "algebraic-geometry", "order_index": 5, "created_at": "2026-07-31T15:26:25.671Z"}, "status": "open"}
% FIRST_POSTED: 2026-09-25
% ENTRY_KIND: statement_only
% !TeX program = lualatex
% Definition and sources only; this file is not an LLM solution attempt.
\documentclass[11pt]{article}
\usepackage[margin=25mm]{geometry}
\usepackage{fontspec}
\setmainfont{Latin Modern Roman}
\usepackage{amsmath,amssymb,mathtools}
\usepackage{unicode-math}
\setmathfont{Latin Modern Math}
\usepackage{xurl,hyperref}
\hypersetup{colorlinks=true,urlcolor=blue}
\setcounter{secnumdepth}{0}
\setlength{\parindent}{0pt}
\setlength{\parskip}{5pt}
\setlength{\emergencystretch}{3em}
\begin{document}
\section*{\texorpdfstring{Kimura–O'Sullivan Finite-Dimensionality Conjecture}{Kimura–O'Sullivan Finite-Dimensionality Conjecture}}\label{30006782}
\subsection{Definitions and mathematical statement}
In the rational Chow-motive category $\mathrm{CHM}(k)_{{\mathbb{Q}}}$, morphisms are rational correspondences modulo rational equivalence, idempotents are split, and tensor product uses products of varieties. Exterior and symmetric powers are the images of the alternating and averaging idempotents of symmetric groups under the ordinary symmetry constraint. The conjecture asserts
\[
 h(X)\cong M_+\oplus M_-,\qquad
 \Lambda^n M_+=0,\qquad\operatorname{Sym}^m M_-=0
\]
for every smooth projective $X/k$, for some $n,m\ge1$. Either summand may be zero. This categorical finite-dimensionality is not simply finiteness of a cohomology vector space. The coefficients are rational but the base field has arbitrary characteristic; no uniform exponents or effective splitting are requested.
\par\smallskip\textit{Formulation and concept references:} \href{https://www.numdam.org/item/SB_2003-2004__46__115_0.pdf}{[S1]}.

\subsection{Short English statement}
Can every smooth projective variety's rational motive be split into an even part with a vanishing exterior power and an odd part with a vanishing symmetric power?
\subsection{Sources}
\begin{itemize}
\item[S1]Motifs de dimension finie [d'après S.-I. Kimura, P. O'Sullivan, ...]; Bourbaki929,March2004; Astérisque299(2005),115–145.
\url{https://www.numdam.org/item/SB_2003-2004__46__115_0.pdf}.
\textit{Formulation source.}
\end{itemize}
\end{document}
